% Generated by Listas/make_statements.py from pratica-lecture-05.md, item 9. Do not edit.
\textbf{Source.} Cabral (2017), \emph{Introduction to Industrial Organization}, 2nd ed., Ch. 7 Challenging exercises, Exercise 7.12 ``Ad games'', p.~181 {[}PDF 193{]}. Answer in the source: none in the book. Reproduced verbatim.



\begin{quote}

\textbf{7.12. Ad games.} Two firms must simultaneously choose their advertising budget; their options are \(H\) or \(L\). Payoffs are as follows: if both choose \(H\), then each gets 5; if both choose \(L\), then each gets 4; if firm 1 chooses \(H\) and firm 2 chooses \(L\), then firm 1 gets 8 and firm 2 gets 1; conversely, if firm 2 chooses \(H\) and firm 1 chooses \(L\), then firm 2 gets 8 and firm 1 gets 1.



\begin{enumerate}

\def\labelenumi{(\alph{enumi})}

\item

  Determine the Nash equilibria of the one-shot game.

\item

  Suppose the game is indefinitely repeated and that the relevant discount factor is \(\delta = .8\). Determine the optimal symmetric equilibrium.

\item

  (challenge question) Now suppose that, for the first 10 periods, firm payoffs are twice the values represented in the above table. What is the optimal symmetric equilibrium?

\end{enumerate}

\end{quote}



\textbf{Note (not part of the exercise).} With the printed payoffs \((H,H)\) gives 5 each and \((L,L)\) only 4 each, so \(H\) is dominant and also efficient. The solution works the exercise as printed.

